\chapter{Environnement de travail et ressources technologiques}

\section*{Introduction}
La réalisation d'un projet informatique nécessite un environnement de travail adéquat, composé de ressources matérielles performantes et d'outils logiciels adaptés. Ce chapitre détaille l'infrastructure technique utilisée durant la période de stage pour mener à bien les phases de conception, développement et rédaction.

\section{Environnement Matériel}
Pour le développement de l'application et la rédaction du rapport, j'ai utilisé un ordinateur portable personnel dont les caractéristiques garantissent une fluidité d'exécution des différents outils.

\textbf{Caractéristiques de l'ordinateur de développement :}
\begin{itemize}
    \item \textbf{Processeur :} Processeurs multicœurs (Intel i7/AMD Ryzen 7 ou équivalents) pour une exécution rapide des compilations et des tests.
    \item \textbf{Mémoire RAM :} 16 Go de RAM pour assurer une fluidité lors de l'exécution simultanée des serveurs locaux, des IDE, et des navigateurs pour le test.
    \item \textbf{Stockage :} Disque SSD (1 To ou plus) pour un accès rapide aux fichiers et aux données.
\end{itemize}

\section{Environnement Logiciel}
Dans cette partie, nous allons présenter les différents outils informatiques utilisés pour le développement, la modélisation et le déploiement.

\subsection{Développement et Tests Locaux}

\noindent
\begin{minipage}[c]{0.3\textwidth}
    \centering
    \includegraphics[width=0.6\textwidth]{images/vscode.png}
    \captionof{figure}{Visual Studio Code}
\end{minipage}
\hfill
\begin{minipage}[c]{0.65\textwidth}
    \textbf{Visual Studio Code (VS Code) :} L'éditeur de code principal utilisé pour ce projet. Sa polyvalence et ses extensions (PHP IntelliSense, GitLink) en font un outil incontournable.
\end{minipage}

\vspace{0.6cm}

\noindent
\begin{minipage}[c]{0.3\textwidth}
    \centering
    \includegraphics[width=0.6\textwidth]{images/xampp.png}
    \captionof{figure}{XAMPP}
\end{minipage}
\hfill
\begin{minipage}[c]{0.65\textwidth}
    \textbf{XAMPP :} Utilisé comme environnement de serveur local (Apache + MySQL + PHP) pour tester l'application avant son déploiement.
\end{minipage}

\vspace{0.6cm}

\subsection{Langages et Technologies Web}

\noindent
\begin{minipage}[c]{0.3\textwidth}
    \centering
    \includegraphics[width=0.5\textwidth]{images/html.png}
    \captionof{figure}{HTML5}
\end{minipage}
\hfill
\begin{minipage}[c]{0.65\textwidth}
    \textbf{HTML5 :} Langage de balisage utilisé pour la structuration sémantique et le contenu des pages web. Il constitue le squelette de l'application.
\end{minipage}

\vspace{0.6cm}

\noindent
\begin{minipage}[c]{0.3\textwidth}
    \centering
    \includegraphics[width=0.5\textwidth]{images/css.png}
    \captionof{figure}{CSS3}
\end{minipage}
\hfill
\begin{minipage}[c]{0.65\textwidth}
    \textbf{CSS3 :} Langage de feuilles de style utilisé pour la mise en forme graphique, le design responsive et l'esthétique globale de l'interface utilisateur.
\end{minipage}

\vspace{0.6cm}

\noindent
\begin{minipage}[c]{0.3\textwidth}
    \centering
    \includegraphics[width=0.5\textwidth]{images/php.png}
    \captionof{figure}{PHP}
\end{minipage}
\hfill
\begin{minipage}[c]{0.65\textwidth}
    \textbf{PHP :} Langage de script côté serveur, utilisé pour le développement de la logique métier, la gestion des sessions utilisateurs et les interactions avec la base de données. Il constitue le moteur dynamique de l'application.
\end{minipage}

\vspace{0.6cm}

\noindent
\begin{minipage}[c]{0.3\textwidth}
    \centering
    \includegraphics[width=0.5\textwidth]{images/js.png}
    \captionof{figure}{JavaScript}
\end{minipage}
\hfill
\begin{minipage}[c]{0.65\textwidth}
    \textbf{JavaScript :} Langage de programmation côté client permettant d'enrichir l'expérience utilisateur via des interactions dynamiques, la validation de formulaires en temps réel et la gestion des menus interactifs.
\end{minipage}

\vspace{0.6cm}

\subsection{Système de Gestion de Base de Données}

\noindent
\begin{minipage}[c]{0.3\textwidth}
    \centering
    \includegraphics[width=0.6\textwidth]{images/mysql.png.jpg}
    \captionof{figure}{MySQL}
\end{minipage}
\hfill
\begin{minipage}[c]{0.65\textwidth}
    \textbf{MySQL :} Système de gestion de base de données relationnelle (SGBDR) robuste et performant. Il assure le stockage structuré, la persistance et l'intégrité des données du système d'information.
\end{minipage}

\vspace{0.6cm}

\subsection{Bibliothèques et Outils Tiers}

\noindent
\begin{minipage}[c]{0.3\textwidth}
    \centering
    \includegraphics[width=0.5\textwidth]{images/bootstrap.png}
    \captionof{figure}{Bootstrap}
\end{minipage}
\hfill
\begin{minipage}[c]{0.65\textwidth}
    \textbf{Bootstrap :} Framework CSS populaire utilisé pour accélérer le développement front-end et garantir une interface parfaitement responsive, s'adaptant à tous les types d'écrans (mobiles, tablettes, desktop).
\end{minipage}

\vspace{0.6cm}

\noindent
\begin{minipage}[c]{0.3\textwidth}
    \centering
    \includegraphics[width=0.5\textwidth]{images/fontawesome.png}
    \captionof{figure}{Font Awesome}
\end{minipage}
\hfill
\begin{minipage}[c]{0.65\textwidth}
    \textbf{Font Awesome :} Bibliothèque d'icônes vectorielles offrant une large gamme de symboles graphiques pour améliorer la lisibilité et l'ergonomie des interfaces utilisateurs.
\end{minipage}

\vspace{0.6cm}

\noindent
\begin{minipage}[c]{0.3\textwidth}
    \centering
    \includegraphics[width=0.6\textwidth]{images/phpmailer.png}
    \captionof{figure}{PHPMailer}
\end{minipage}
\hfill
\begin{minipage}[c]{0.65\textwidth}
    \textbf{PHPMailer :} Bibliothèque PHP complète et sécurisée pour la gestion de l'envoi d'emails. Elle est utilisée pour les notifications automatiques et l'envoi de confirmations aux utilisateurs.
\end{minipage}

\subsection{Déploiement et Hébergement}

\noindent
\begin{minipage}[c]{0.3\textwidth}
    \centering
    \includegraphics[width=0.6\textwidth]{images/ovh.jpg}
    \captionof{figure}{OVHcloud}
\end{minipage}
\hfill
\begin{minipage}[c]{0.65\textwidth}
    \textbf{OVHcloud :} Fournisseur d'hébergement choisi pour la mise en ligne de l'application. Sa fiabilité et ses serveurs sécurisés garantissent la disponibilité de la plateforme.
\end{minipage}

\section*{Conclusion}
Ce chapitre a permis de dresser le panorama technique du projet. De l'environnement matériel robuste à la suite logicielle intégrant des technologies web fondamentales (HTML, CSS, JavaScript et MySQL), l'infrastructure a été pensée pour allier performance, sécurité et maintenabilité.
